\documentclass[10pt,final,hidelinks]{article}
\usepackage{lmodern}
\usepackage[T1]{fontenc}
\usepackage[english]{babel}
\usepackage[utf8]{inputenc}
\usepackage[activate={true,nocompatibility},final,tracking=true,kerning=true,
    spacing=true,factor=1100,stretch=10,shrink=10]{microtype}
\microtypecontext{spacing=nonfrench}
\usepackage[margin=.5in]{geometry}
\usepackage{graphicx}
\graphicspath{{img/}} 
\usepackage[noend]{algorithm2e}
\usepackage{caption}
\usepackage{subcaption}
\usepackage{booktabs}
\usepackage{array}
\usepackage{mathtools}
\usepackage{subfiles}
\usepackage{enumitem}
\usepackage{commath}
\usepackage{appendix}
\renewcommand\appendixtocname{Appendices}
\usepackage{amsmath}
\usepackage{amsthm}
\usepackage{amssymb}
%\usepackage[intoc,english]{nomencl}
%\makenomenclature
\usepackage{pgfplots}
\usepackage{tabu} 
\usepackage{tikz}
\usepackage{tikzscale}
\usepackage[square,numbers]{natbib}
\bibliographystyle{plain}
\usepackage{siunitx}
\numberwithin{equation}{section}
\usepackage{hyperref}
\usepackage{cleveref}

\pgfplotsset{compat=newest}

\newcommand{\COMMENT}[1]{}

\usepackage{functionnotation}
\usepackage{algorithmnotation}
\usepackage{matrixnotation}
\usepackage[fancy]{setnotation}
\usepackage[fancy]{relationnotation}
\usepackage{approxsetnotation}
\usepackage{approxrelationnotation}
\usepackage[section]{envnotation}

\hypersetup{
    pdftitle={Bernoulli sets and their corresponding random binary classification measures.},
    pdfauthor={Alexander Towell},               % author
    pdfsubject={computer science},              % subject of the %document
    pdfkeywords={
        probabilistic data structure,
        abstract data type,
        approximate set,
        bloom filter,
        perfect hash function,
        perfect hash filter},                   % keywords
    colorlinks=true,                            % false: boxed links;
    linkcolor=magenta,
    citecolor=green,                            % color of links to 
    filecolor=blue,                             % color of file links
    urlcolor=green                              % color of external
}

\title
{
    Bernoulli sets: Probabilistic set approximations with controlled error rates\\
    \large Theory, algebra, and applications to information retrieval
}
\author
{
    Alexander Towell\\
    \texttt{atowell@siue.edu}
}
\date{}

\begin{document}
\maketitle
\begin{abstract}
We present a comprehensive theory of Bernoulli sets---probabilistic approximations of classical sets with controlled false positive and false negative rates.
We establish the mathematical foundations including the probability spaces, confusion matrix structures, and information-theoretic properties that govern these approximations.
We derive the error propagation laws for Boolean algebra operations (union, intersection, complement) and characterize the resulting higher-order approximations.
We analyze the distribution of binary classification measures (precision, recall) when Bernoulli sets are used in practice.
Finally, we demonstrate a concrete application to Boolean search in information retrieval systems, showing how Bernoulli sets enable space-efficient indexing with quantifiable accuracy guarantees and natural privacy properties.
Our framework unifies existing probabilistic data structures like Bloom filters within a rigorous mathematical foundation while revealing new insights about rank, inferability, and composition.
\end{abstract}

%\microtypesetup{protrusion=false}
\tableofcontents
%\microtypesetup{protrusion=true}
%\listoftables
%\listoffigures

%\include{nom}
%\printnomenclature

\subfile{sections/intro}
\subfile{sections/algebra_of_sets}
%\subfile{sections/aset_model}
\subfile{sections/bernoulli_model}
\subfile{sections/confusion_matrix_analysis}
\subfile{sections/first_order_model}
\subfile{sections/second_order_pos_neg_model}
\subfile{sections/derived_distributions}
\subfile{sections/derived_distributions_higher_order}
\subfile{sections/relational}
\subfile{sections/interval}
\subfile{sections/aset_adt}
\subfile{sections/application_boolean_search}
\subfile{sections/appendix}

\bibliography{references}

\end{document}
